% !TeX root = ../main.tex

\chapter{Representation Theorem}

\section{A New Expansion in Centered Blocks}
% =================================================

Consider an ordinary formal power series

\[
f(x)
=
f(0)
+
t_1x
+
t_2x^2
+
t_3x^3
+\cdots.
\]

We can ask whether it is possible to write

\[
\boxed{
f(x)
=
f(0)
+
\sum_{n\geq1}
c_n
\left(
e^{\lambda x^n}-1
\right).
}
\]

Expanding the right-hand side gives

\[
\sum_{n\geq1}
c_n
\sum_{k\geq1}
\frac{\lambda^k}{k!}x^{nk}.
\]

Therefore the coefficient \(t_m\) must satisfy

\[
\boxed{
t_m
=
\sum_{n\mid m}
c_n
\frac{\lambda^{m/n}}{(m/n)!}.
}
\]

Among the divisors of \(m\) is \(n=m\).

The corresponding contribution is

\[
c_m\lambda.
\]

Hence

\[
t_m
=
\lambda c_m
+
\sum_{\substack{n\mid m\\n<m}}
c_n
\frac{\lambda^{m/n}}{(m/n)!}.
\]

Solving for \(c_m\),

\[
\boxed{
c_m
=
\frac{1}{\lambda}
\left[
t_m
-
\sum_{\substack{n\mid m\\n<m}}
c_n
\frac{\lambda^{m/n}}{(m/n)!}
\right].
}
\]

This gives a recursive procedure for constructing the coefficients
\(c_m\).

% =================================================
\section{Formal Representation Theorem}
% =================================================

The previous recursion gives an interesting result.

\begin{theorem}

Let

\[
\lambda\neq0.
\]

Every formal power series

\[
f(x)
=
f(0)+\sum_{m\geq1}t_mx^m
\]

has a unique formal representation of the form

\[
\boxed{
f(x)
=
f(0)
+
\sum_{n\geq1}
c_n
\left(
e^{\lambda x^n}-1
\right).
}
\]

\end{theorem}

\begin{proof}

For \(m=1\),

\[
t_1=\lambda c_1,
\]

so

\[
c_1=\frac{t_1}{\lambda}.
\]

Suppose

\[
c_1,\ldots,c_{m-1}
\]

have already been determined.

Then

\[
t_m
=
\lambda c_m
+
\sum_{\substack{n\mid m\\n<m}}
c_n
\frac{\lambda^{m/n}}{(m/n)!}.
\]

Every term in the sum is already known.

Since

\[
\lambda\neq0,
\]

we can solve uniquely for \(c_m\).

Proceeding recursively determines every coefficient \(c_m\).

\end{proof}

\begin{remark}

This theorem concerns formal power series.

It does not automatically imply that the corresponding infinite
series converges as a function for a particular value of \(x\).

Understanding when the formal representation becomes an analytic
representation is a separate and important problem.

\end{remark}

% =================================================
\section{Prime Indices in the Taylor Transform}
% =================================================

The divisor formula becomes especially simple for prime indices.

Suppose \(p\) is prime.

Its only positive divisors are

\[
1
\qquad\text{and}\qquad
p.
\]

Therefore

\[
t_p
=
c_1\frac{\lambda^p}{p!}
+
c_p\lambda.
\]

Hence

\[
\boxed{
c_p
=
\frac{1}{\lambda}
\left(
t_p
-
c_1\frac{\lambda^p}{p!}
\right).
}
\]

For a composite number, more terms appear.

For example, the divisors of \(6\) are

\[
1,2,3,6.
\]

Therefore

\[
t_6
=
c_1\frac{\lambda^6}{6!}
+
c_2\frac{\lambda^3}{3!}
+
c_3\frac{\lambda^2}{2!}
+
c_6\lambda.
\]

Thus the arithmetic complexity of the index controls the number of
interactions between coefficients.

% =================================================
\section{Linear Independence}
% =================================================

The centered blocks are also linearly independent.

\begin{proposition}

For \(\lambda\neq0\), the functions

\[
\psi_n(x)
=
e^{\lambda x^n}-1,
\qquad n\geq1,
\]

are linearly independent with respect to finite linear combinations.

\end{proposition}

\begin{proof}

Suppose

\[
c_1\psi_1(x)
+
c_2\psi_2(x)
+\cdots+
c_N\psi_N(x)
=
0.
\]

Assume not all coefficients are zero.

Let \(m\) be the smallest index for which

\[
c_m\neq0.
\]

The lowest power of \(x\) contributed by \(\psi_m\) is

\[
\lambda x^m.
\]

Blocks with index greater than \(m\) cannot contribute to \(x^m\).

Blocks with index smaller than \(m\) have coefficient zero by the
minimality of \(m\).

Therefore the coefficient of \(x^m\) in the entire expression is

\[
\lambda c_m.
\]

Since

\[
\lambda\neq0
\qquad\text{and}\qquad
c_m\neq0,
\]

this coefficient cannot vanish.

This is a contradiction.

Hence every \(c_n\) must be zero.

\end{proof}

% =================================================
